% !TeX program = xelatex
\documentclass[UTF8,a4paper,zihao=-4]{ctexart}

\usepackage{geometry}
\usepackage{setspace}
\usepackage{longtable}
\usepackage{booktabs}
\usepackage{array}
\usepackage{xcolor}
\usepackage{enumitem}
\geometry{left=2.6cm,right=2.6cm,top=2.5cm,bottom=2.5cm}
\setstretch{1.18}
\pagestyle{plain}
\setlength{\parindent}{2em}
\setlist[itemize]{leftmargin=2em,itemsep=0.2em,topsep=0.2em}

\begin{document}

\begin{center}
{\zihao{3}\heiti 人工智能工具使用详情}
\end{center}

\section*{一、使用工具说明}

本文在写作整理过程中使用了人工智能工具，主要用于文字梳理、表达润色、结构检查和部分说明性段落的初稿组织。具体信息如下。

\begin{longtable}{p{3.1cm}p{10.2cm}}
\toprule
项目 & 内容 \\
\midrule
工具名称 & OpenAI Codex \\
版本/模型 & GPT-5 Codex \\
开发机构/公司 & OpenAI \\
使用日期 & 2026-06-01 \\
主要用途 & 辅助整理论文文字、检查章节衔接、生成或改写部分说明性表述；不用于直接生成原始数据、计算结果和核心模型参数。 \\
\bottomrule
\end{longtable}

说明：人工智能工具输出内容均经过人工核对、删改和重新组织。论文中的数据处理、指标数值、模型选择、图表结果和结论判断均以本文实际建模结果为准。

\section*{二、具体使用目的和环节}

结合正文中已标注人工智能工具引用的位置，主要使用环节如下。

\begin{longtable}{p{1.0cm}p{3.2cm}p{8.7cm}}
\toprule
序号 & 使用位置 & 使用目的 \\
\midrule
1 & 研究背景末段 & 根据烧结过程数据特点，辅助整理“数据驱动预测模型和优化调控模型”的背景表达，使其与后文预测、调控任务衔接更自然。 \\
2 & 问题二时滞分析 & 辅助把风箱位置差异、物料移动、烟气汇集和时间配准之间的逻辑关系表述清楚。 \\
3 & 时滞配准结果解释 & 辅助整理“配准效果从相关性、工艺认识、模型指标三方面判断”的说明文字。 \\
4 & 异常样本处理 & 辅助表达极端 CO 样本可能来源于非稳态工况或检测波动，以及为什么主模型采用剔除异常后的样本。 \\
5 & 模型推广部分 & 辅助整理模型推广到其他多测点、多变量、明显时滞工业过程的表述。 \\
\bottomrule
\end{longtable}

\section*{三、关键交互记录摘录}

以下为关键交互记录的整理摘录。由于实际写作中存在多轮复制、粘贴、改写和删改，记录中保留部分零散表达，不完全按正式论文格式呈现。

\subsection*{记录 1：研究背景和任务衔接}

\noindent\textbf{用户提示：}

“烧结烟气 CO 排放受机速、风箱负压、温度场和大烟道状态等多类因素共同影响，问题一二不是孤立的，要服务问题三。背景最后要自然一点，别太像 AI。”

\noindent\textbf{AI 回复要点：}

可把背景最后落到“已有生产数据为预测和调控提供基础”，再说风箱沿烧结机方向分布，影响具有传递性和滞后性，因此需要分析风箱负压、温度和 CO 排放之间的关系。不要直接写成很空的“构建状态表征--排放预测--负压优化”。

\noindent\textbf{采纳情况：}

采纳了“数据驱动预测与优化调控衔接”的思路，但删去了较模板化的表达，并根据正文已有结果改写为现在的研究背景末段。

\subsection*{记录 2：问题二时滞解释}

\noindent\textbf{用户提示：}

“问题二这里要说清楚风箱位置和 CO 预测之间的关系，不要只写建立模型。前面的风箱、后面的风箱，和检测位置有时间差。”

\noindent\textbf{AI 回复要点：}

可以写：18 个风箱沿烧结机长度方向依次布置，各风箱距离外排烟气检测位置不同，物料移动、燃烧带推进和烟气汇集都需要时间。靠前风箱的影响经过较长传递过程，靠后风箱影响相对更快。因此需要构造不同滞后步数下的特征。

\noindent\textbf{采纳情况：}

基本采纳了这段逻辑，但人工调整了文字顺序，使其与“问题二分析”和后文公式部分一致。

\subsection*{记录 3：时滞配准结果怎么解释}

\noindent\textbf{用户提示：}

“这个表出来以后要解释，不能只摆表。需要说为什么前段中段后段的滞后步数不一样，还有配准效果怎么看。”

\noindent\textbf{AI 回复要点：}

建议从三个方面写：一是配准后变量与 CO 浓度相关性是否增强；二是不同风箱的滞后步数是否符合前段较长、后段较短的工艺认识；三是输入预测模型后 MAE、RMSE、$R^2$ 是否改善。可补一句“最终得到较合理的时间对应关系”。

\noindent\textbf{采纳情况：}

采纳了三点判断标准。人工将“相关性是否增强”和“模型指标是否优于未配准模型”保留为评价口径，并结合表格结果进行了压缩。

\subsection*{记录 4：异常 CO 样本处理}

\noindent\textbf{用户提示：}

“CO 浓度极端低值和极端高值可能是异常，也可能是非稳态。这里不要写得太绝对。主模型剔除异常，但要保留对比分析。”

\noindent\textbf{AI 回复要点：}

可以写“异常候选点”，避免直接认定为错误数据。调控模型优先采用 $I_t=0$ 的样本训练，剔除前后的模型结果放入稳健性分析。这样处理不是否定极端样本存在，而是将其作为可能的非稳态或检测异常单独讨论。

\noindent\textbf{采纳情况：}

采纳“异常候选点”和“不是否定极端样本”的表述。人工补充了阈值定义和稳健性分析结果。

\subsection*{记录 5：模型推广部分}

\noindent\textbf{用户提示：}

“模型评价这里参考一等奖论文，写优点、不足、推广。推广不要太虚，但是可以说能推广到类似工业过程。”

\noindent\textbf{AI 回复要点：}

可按数据层面、模型层面、约束层面、应用层面写。应用层面可写：适用于多测点、多变量、明显时滞的工业过程，如烧结终点控制、烟气 NOx 排放预测、球团焙烧过程优化等。前提是有连续运行数据、明确可调变量和目标指标。

\noindent\textbf{采纳情况：}

采纳了分类思路和部分例子，人工删改了过于泛化的句子，使推广部分更贴合本文数据和模型。

\section*{四、人工采纳和修改情况}

人工智能工具生成的内容没有直接作为最终论文结论使用。具体处理包括：

\begin{itemize}
  \item 对背景、问题分析、模型评价等文字进行了人工删改，删除了过于模板化、空泛或重复的表达；
  \item 对所有数值结果、表格数据、模型指标、优化结果进行人工核对，未使用人工智能工具生成数值；
  \item 对参考文献、符号说明、摘要和正文引用位置进行了人工检查；
  \item 对 AI 建议的部分标题和句式进行了调整，使其符合全文写作风格；
  \item 对“异常值”“非稳态”“检测异常”等表述保留不确定性，避免把模型假设写成事实判断。
\end{itemize}

综上，人工智能工具主要承担文字辅助和结构整理作用，最终内容由作者根据赛题要求、建模结果和论文整体逻辑进行人工判断、选择和修改。

\end{document}
